\subsection{Smith 亏损轮廓与局部 zeta 完备不变量}\label{subsec:conclusion-smith-bounded-profile-local-zeta}
沿用定理 \ref{thm:xi-smith-loss-spectrum-conjugate-partition-realization}、定理 \ref{thm:xi-kernel-loss-divisibility-valuation-rational-euler} 与定理 \ref{thm:xi-smith-local-euler-complete-invariant} 的记号。前述链条已经给出：固定长度 Smith 亏损谱的共轭分拆恢复、模核 Dirichlet 级数的有理 Euler 局部因子，以及局部 Euler 因子对 Smith 正规形的完备性。由此还可进一步抽出一组彼此锁定的后继结果：固定行秩口径下亏损轮廓的完全刻画、估值分拆与差分剖面的 Young 对偶、归一化模核 Dirichlet 级数的正系数有限 Euler 修正与平均阶主项、归一化核大小 Dirichlet 级数与 torsion 余核的完全等价、核大小 zeta 的自由极与有限 Euler 变形、坏素数局部因子对全部 Smith 不变量的可逆恢复、坏素集合的亏损谱刻画，以及 prime-power Ramanujan 阴影对局部 zeta 与有限 prime-stage 截断不完备性的精确重写。

\begin{theorem}[固定行秩口径下 Smith 信息亏损轮廓的完全刻画]\label{thm:conclusion-smith-bounded-rank-loss-profile-characterization}
\leanverified{paper\_conclusion\_smith\_bounded\_rank\_loss\_profile\_characterization}
固定整数 \(m\ge 1\)。设
\[
f:\ZZ_{\ge 0}\to \ZZ_{\ge 0},
\qquad
\Delta(k):=f(k)-f(k-1)\quad(k\ge 1).
\]
则下列条件等价：
\begin{enumerate}
  \item 存在某个满行秩整数矩阵 \(A\in M_{m\times n}(\ZZ)\) 与某个素数 \(p\)，使得
  \[
  f(k)=f_p(k):=\sum_{i=1}^{m}\min(k,e_i),
  \qquad
  e_i:=v_p(d_i),
  \]
  其中 \(d_1\mid\cdots\mid d_m\) 为 \(A\) 的 Smith 不变量；
  \item \(f(0)=0\)，并且 \(\Delta(k)\in\{0,1,\dots,m\}\) 对所有 \(k\ge 1\) 成立，且 \((\Delta(k))_{k\ge 1}\) 单调不增并最终为零。
\end{enumerate}
在这些条件成立时，长度恰为 \(m\) 的非负估值序列
\[
e_1^\downarrow\ge\cdots\ge e_m^\downarrow\ge 0
\]
由 \(f\) 唯一决定。
\end{theorem}
\begin{proof}
若 \(f=f_p\)，则由定理 \ref{thm:killo-smith-loss-spectrum}，
\[
\Delta(k)=f_p(k)-f_p(k-1)=\#\{i:\ e_i\ge k\}.
\]
故 \(\Delta(k)\) 取值于 \(\{0,\dots,m\}\)，单调不增，且当 \(k>\max_i e_i\) 时恒为 \(0\)。因此 \(1\Rightarrow 2\)。

反之，若 \(2\) 成立，则可定义
\[
e_j^\downarrow:=\#\{k\ge 1:\ \Delta(k)\ge j\}
\qquad(1\le j\le m).
\]
由定理 \ref{thm:xi-smith-loss-spectrum-conjugate-partition-realization}，
\[
f(k)=\sum_{j=1}^{m}\min(k,e_j^\downarrow)
\qquad(k\ge 0).
\]
固定任意素数 \(p\)，取
\[
A:=\operatorname{diag}(p^{e_1^\downarrow},\dots,p^{e_m^\downarrow})\in M_m(\ZZ),
\]
则 \(A\) 满秩且其 \(p\)-主 Smith 估值恰为 \((e_1^\downarrow,\dots,e_m^\downarrow)\)，从而 \(f=f_p\)。唯一性同样由定理 \ref{thm:xi-smith-loss-spectrum-conjugate-partition-realization} 给出。证毕。
\end{proof}

\begin{corollary}[信息亏损轮廓与估值分拆的 Young 对偶]\label{cor:conclusion-smith-loss-profile-young-duality}
在定理 \ref{thm:conclusion-smith-bounded-rank-loss-profile-characterization} 的设定下，设
\[
e_1^\downarrow\ge\cdots\ge e_m^\downarrow\ge 0
\]
为按降序排列后的 \(p\)-主 Smith 估值，且
\[
\Delta_p(k):=f_p(k)-f_p(k-1).
\]
则对每个 \(1\le j\le m\) 都有
\[
\boxed{\;
e_j^\downarrow=\#\{k\ge 1:\ \Delta_p(k)\ge j\}.\;}
\]
更一般地，对任意 \(1\le \ell\le m\)，有
\[
\boxed{\;
\sum_{j=1}^{\ell}e_j^\downarrow
=
\sum_{k\ge 1}\min(\ell,\Delta_p(k)).\;}
\]
\end{corollary}
\begin{proof}
考虑 Ferrers 点阵
\[
\mathcal Y_p:=\{(j,k)\in\{1,\dots,m\}\times \NN:\ e_j^\downarrow\ge k\}.
\]
其第 \(k\) 列高度正是
\[
\#\{j:\ e_j^\downarrow\ge k\}=\Delta_p(k),
\]
而第 \(j\) 行长度正是 \(e_j^\downarrow\)。故
\[
e_j^\downarrow=\#\{k\ge 1:\ \Delta_p(k)\ge j\}.
\]
再对前 \(\ell\) 行计数，即得
\[
\sum_{j=1}^{\ell}e_j^\downarrow
=
\sum_{k\ge 1}\#\{1\le j\le \ell:\ e_j^\downarrow\ge k\}
=
\sum_{k\ge 1}\min(\ell,\Delta_p(k)).
\]
证毕。
\end{proof}

\begin{theorem}[核大小 Dirichlet 级数的 Euler 分解与有理局部因子]\label{thm:conclusion-smith-kernel-dirichlet-euler-rational-local-factors}
设 \(A\in M_{m\times n}(\ZZ)\) 满足 \(\operatorname{rank}_{\QQ}(A)=m\)，其 Smith 不变量为
\[
d_1\mid\cdots\mid d_m.
\]
定义
\[
\mathcal Z_A(s):=\sum_{b\ge 1}\frac{\#\ker(A_b)}{b^s},
\qquad
A_b:(\ZZ/b\ZZ)^n\to(\ZZ/b\ZZ)^m
\qquad
(\Re s>n-m+1).
\]
对每个素数 \(p\)，记
\[
e_i(p):=v_p(d_i),
\qquad
f_p(k):=\sum_{i=1}^{m}\min(k,e_i(p)),
\]
\[
E_p^{\max}:=\max_i e_i(p),
\qquad
E_p^{\Sigma}:=\sum_{i=1}^{m}e_i(p).
\]
则 \(\mathcal Z_A\) 具有 Euler 乘积
\[
\boxed{\;
\mathcal Z_A(s)=\prod_p \mathcal Z_{A,p}(s),\qquad
\mathcal Z_{A,p}(s)=\sum_{k\ge 0}p^{-k(s-n+m)+f_p(k)}.\;}
\]
并且每个局部因子都是 \(p^{-s}\) 的有理函数，更精确地，
\[
\boxed{\;
\mathcal Z_{A,p}(s)
=
\sum_{k=0}^{E_p^{\max}-1}p^{-k(s-n+m)+f_p(k)}
\;+\;
\frac{p^{-E_p^{\max}(s-n+m)+E_p^{\Sigma}}}{1-p^{-(s-n+m)}}.\;}
\]
\end{theorem}
\begin{proof}
由推论 \ref{cor:killo-kernel-size-general-modulus}，
\[
\#\ker(A_b)=b^{\,n-m}\prod_{i=1}^{m}\gcd(d_i,b).
\]
右端关于 \(b\) 是乘法函数，故 \(\mathcal Z_A\) 具有 Euler 分解。

对 \(b=p^k\)，上式化为
\[
\#\ker(A_{p^k})
=
p^{k(n-m)}\prod_{i=1}^{m}p^{\min(k,e_i(p))}
=
p^{k(n-m)+f_p(k)},
\]
故
\[
\mathcal Z_{A,p}(s)=\sum_{k\ge 0}p^{-ks+k(n-m)+f_p(k)}
=
\sum_{k\ge 0}p^{-k(s-n+m)+f_p(k)}.
\]
当 \(k\ge E_p^{\max}\) 时，\(\min(k,e_i(p))=e_i(p)\) 对所有 \(i\) 都成立，因此
\[
f_p(k)=E_p^{\Sigma}.
\]
于是尾部成为公比 \(p^{-(s-n+m)}\) 的几何级数：
\[
\sum_{k\ge E_p^{\max}}p^{-k(s-n+m)+E_p^{\Sigma}}
=
\frac{p^{-E_p^{\max}(s-n+m)+E_p^{\Sigma}}}{1-p^{-(s-n+m)}}.
\]
证毕。
\end{proof}

\begin{corollary}[局部 zeta 因子是 \texorpdfstring{$p$}{p}-主 Smith 谱的完备不变量]\label{cor:conclusion-smith-local-zeta-complete-invariant}
在定理 \ref{thm:conclusion-smith-kernel-dirichlet-euler-rational-local-factors} 的设定下，对每个素数 \(p\) 定义形式幂级数
\[
\mathfrak Z_{A,p}(X):=\sum_{k\ge 0}p^{f_p(k)}X^k.
\]
则 \(\mathfrak Z_{A,p}(X)\) 唯一决定 \(A\) 的 \(p\)-主 Smith 估值多重集
\[
\{e_i(p)\}_{i=1}^{m}.
\]
因而有限个坏素数上的局部因子族
\[
\{\mathfrak Z_{A,p}(X):\ p\mid d_1\cdots d_m\}
\]
已经足以可逆恢复全部 Smith 不变量。
\end{corollary}
\begin{proof}
对每个 \(k\ge 0\)，\(\mathfrak Z_{A,p}(X)\) 的 \(X^k\) 项系数恰为 \(p^{f_p(k)}\)，故
\[
f_p(k)=\log_p\!\bigl([X^k]\mathfrak Z_{A,p}(X)\bigr).
\]
于是整个局部亏损轮廓 \((f_p(k))_{k\ge 0}\) 都可逐项恢复。再由定理 \ref{thm:conclusion-smith-bounded-rank-loss-profile-characterization}，该轮廓唯一决定长度为 \(m\) 的估值序列 \((e_i(p))_{i=1}^{m}\)。

另一方面，仅有有限多个素数整除 \(d_1\cdots d_m\)。对其余素数，\(e_i(p)=0\) 对所有 \(i\) 都成立，故局部因子恒为平凡因子。把全部坏素数上的估值数据逐素数拼合，便恢复全部 Smith 不变量。证毕。
\end{proof}

\begin{theorem}[归一化核大小 Dirichlet 刚性与 torsion 余核的完全等价]\label{thm:conclusion-smith-normalized-kernel-dirichlet-torsion-equivalence}
设 \(A,B\in M_{m\times n}(\ZZ)\) 满足
\[
\rank_{\QQ}(A)=\rank_{\QQ}(B)=m.
\]
对每个 \(b\ge 1\)，定义
\[
A_b:(\ZZ/b\ZZ)^n\to(\ZZ/b\ZZ)^m,
\qquad
h_A(b):=b^{-(n-m)}\#\ker(A_b),
\]
\[
B_b:(\ZZ/b\ZZ)^n\to(\ZZ/b\ZZ)^m,
\qquad
h_B(b):=b^{-(n-m)}\#\ker(B_b),
\]
并记
\[
H_A(s):=\sum_{b\ge 1}\frac{h_A(b)}{b^s},
\qquad
H_B(s):=\sum_{b\ge 1}\frac{h_B(b)}{b^s}.
\]
则下列四条等价：
\begin{enumerate}
  \item \(H_A(s)=H_B(s)\) 作为 Dirichlet 级数恒等；
  \item \(h_A(b)=h_B(b)\) 对一切 \(b\ge 1\) 成立；
  \item 对每个素数 \(p\) 与每个 \(k\ge 1\)，有
  \[
  v_p\!\bigl(h_A(p^k)\bigr)=v_p\!\bigl(h_B(p^k)\bigr);
  \]
  \item
  \[
  \operatorname{coker}(A)_{\mathrm{tor}}
  \cong
  \operatorname{coker}(B)_{\mathrm{tor}}.
  \]
\end{enumerate}
换言之，归一化核大小 Dirichlet 级数恰好是满行秩整数矩阵 torsion 余核的完备不变量。
\end{theorem}
\begin{proof}
由 Dirichlet 级数在收敛半平面中的系数唯一性，\(1\Rightarrow 2\) 立即成立，而 \(2\Rightarrow 1\) 显然。

若 \(2\) 成立，则对每个素数 \(p\) 与每个 \(k\ge 1\) 都有
\[
h_A(p^k)=h_B(p^k),
\]
故
\[
v_p\!\bigl(h_A(p^k)\bigr)=v_p\!\bigl(h_B(p^k)\bigr),
\]
即得 \(2\Rightarrow 3\)。

现证 \(3\Rightarrow 4\)。记
\[
e_i^A(p):=v_p(d_i^A),
\qquad
e_i^B(p):=v_p(d_i^B),
\]
其中 \(d_i^A,d_i^B\) 分别为 \(A,B\) 的 Smith 不变量。由定理 \ref{thm:conclusion-smith-kernel-dirichlet-euler-rational-local-factors} 的局部核公式，
\[
\#\ker(A_{p^k})=p^{k(n-m)+f_{A,p}(k)},
\qquad
f_{A,p}(k):=\sum_{i=1}^{m}\min(k,e_i^A(p)),
\]
故
\[
v_p\!\bigl(h_A(p^k)\bigr)=f_{A,p}(k),
\qquad
v_p\!\bigl(h_B(p^k)\bigr)=f_{B,p}(k).
\]
由假设 \(3\) 得 \(f_{A,p}(k)=f_{B,p}(k)\) 对一切 \(p,k\) 成立。于是
\[
\Delta_{A,p}(k):=f_{A,p}(k)-f_{A,p}(k-1)=\#\{i:e_i^A(p)\ge k\},
\]
\[
\Delta_{B,p}(k):=f_{B,p}(k)-f_{B,p}(k-1)=\#\{i:e_i^B(p)\ge k\}
\]
也逐项相等，从而
\[
\#\{i:e_i^A(p)=k\}
=
\Delta_{A,p}(k)-\Delta_{A,p}(k+1)
=
\Delta_{B,p}(k)-\Delta_{B,p}(k+1)
=
\#\{i:e_i^B(p)=k\}.
\]
故每个素数方向上的估值多重集一致，于是 torsion Smith 不变量完全相同，得到
\[
\operatorname{coker}(A)_{\mathrm{tor}}
\cong
\operatorname{coker}(B)_{\mathrm{tor}}.
\]

最后证 \(4\Rightarrow 2\)。若 torsion 余核同构，则 \(A,B\) 具有相同的全部非平凡 Smith 不变量，记为
\[
t_1,\dots,t_r.
\]
由推论 \ref{cor:killo-kernel-size-general-modulus}，
\[
h_A(b)=\prod_{j=1}^{r}\gcd(t_j,b)=h_B(b)
\qquad(\forall b\ge 1).
\]
故 \(2\) 成立。证毕。
\end{proof}

\begin{theorem}[核大小 zeta 的自由极与有限 torsion Euler 变形]\label{thm:conclusion-smith-kernel-zeta-free-pole-finite-euler-deformation}
设 \(A\in M_{m\times n}(\ZZ)\) 满足 \(\rank_{\QQ}(A)=m\)，并定义
\[
Z_A(s):=\sum_{b\ge 1}\frac{\#\ker(A_b)}{b^s}.
\]
记其 Smith 不变量为
\[
d_1\mid\cdots\mid d_m,
\]
并对每个素数 \(p\) 记
\[
e_i(p):=v_p(d_i),
\qquad
f_p(k):=\sum_{i=1}^{m}\min(k,e_i(p)),
\]
\[
E_p^{\max}:=\max_i e_i(p),
\qquad
E_p^{\Sigma}:=\sum_i e_i(p),
\qquad
S_A:=\{p:\exists i,\ e_i(p)>0\}.
\]
则有分解
\[
Z_A(s)
=
\zeta(s-n+m)
\prod_{p\in S_A}
\mathcal E_{A,p}\!\bigl(p^{-(s-n+m)}\bigr),
\]
其中每个局部修正因子都是显式多项式
\[
\mathcal E_{A,p}(X)
=
1+\sum_{k=1}^{E_p^{\max}-1}
\bigl(p^{f_p(k)}-p^{f_p(k-1)}\bigr)X^k
\;
\bigl(p^{E_p^\Sigma}-p^{f_p(E_p^{\max}-1)}\bigr)X^{E_p^{\max}}.
\]
特别地：
\begin{enumerate}
  \item \(Z_A(s)\) 的极点部分完全由自由秩 \(n-m\) 决定，即普适因子 \(\zeta(s-n+m)\)；
  \item 一切 torsion 信息都被压缩到有限个坏素数上的 Euler 修正多项式中；
  \item 商
  \[
  \frac{Z_A(s)}{\zeta(s-n+m)}
  \]
  是一个有限 Euler 变形，并且它作为 Dirichlet 级数完全决定 \(\operatorname{coker}(A)_{\mathrm{tor}}\)。
\end{enumerate}
\end{theorem}
\begin{proof}
由定理 \ref{thm:conclusion-smith-kernel-dirichlet-euler-rational-local-factors}，
\[
Z_A(s)=\prod_p Z_{A,p}(s),
\qquad
Z_{A,p}(s)
=
\sum_{k=0}^{E_p^{\max}-1}p^{-k(s-n+m)+f_p(k)}
\;+\;
\frac{p^{-E_p^{\max}(s-n+m)+E_p^\Sigma}}{1-p^{-(s-n+m)}}.
\]
记
\[
X:=p^{-(s-n+m)}.
\]
则
\[
Z_{A,p}(s)=\sum_{k=0}^{E_p^{\max}-1}p^{f_p(k)}X^k+\frac{p^{E_p^\Sigma}X^{E_p^{\max}}}{1-X}.
\]
乘以 \(1-X\) 可得
\[
(1-X)Z_{A,p}(s)
=
1+\sum_{k=1}^{E_p^{\max}-1}
\bigl(p^{f_p(k)}-p^{f_p(k-1)}\bigr)X^k
\;
\bigl(p^{E_p^\Sigma}-p^{f_p(E_p^{\max}-1)}\bigr)X^{E_p^{\max}},
\]
即
\[
Z_{A,p}(s)=\frac{\mathcal E_{A,p}(X)}{1-X}.
\]

若 \(p\notin S_A\)，则 \(e_i(p)=0\) 对所有 \(i\) 成立，从而 \(E_p^{\max}=E_p^\Sigma=0\) 且 \(Z_{A,p}(s)=(1-X)^{-1}\)。因此
\[
Z_A(s)
=
\prod_p \frac{1}{1-p^{-(s-n+m)}}
\prod_{p\in S_A}\mathcal E_{A,p}\!\bigl(p^{-(s-n+m)}\bigr)
=
\zeta(s-n+m)\prod_{p\in S_A}\mathcal E_{A,p}\!\bigl(p^{-(s-n+m)}\bigr).
\]
这就证明了前两条。

再注意
\[
H_A(s):=\sum_{b\ge 1}\frac{h_A(b)}{b^s}
=
\sum_{b\ge 1}\frac{\#\ker(A_b)}{b^{s+n-m}}
=
Z_A(s+n-m).
\]
故
\[
H_A(s)=\zeta(s)\prod_{p\in S_A}\mathcal E_{A,p}(p^{-s}).
\]
因此有限 Euler 变形
\[
\frac{Z_A(s)}{\zeta(s-n+m)}
\]
与归一化级数 \(H_A\) 彼此唯一决定。再由定理 \ref{thm:conclusion-smith-normalized-kernel-dirichlet-torsion-equivalence}，它作为 Dirichlet 级数完全决定 \(\operatorname{coker}(A)_{\mathrm{tor}}\)。证毕。
\end{proof}

\begin{theorem}[规范化模核 Dirichlet 级数的正系数有限 Euler 修正]\label{thm:conclusion-smith-normalized-kernel-positive-finite-euler-correction}
设 \(A\in M_{m\times n}(\ZZ)\) 满足 \(\operatorname{rank}_{\QQ}(A)=m\)，其 Smith 不变量为
\[
d_1\mid d_2\mid\cdots\mid d_m.
\]
定义
\[
h_A(b):=\frac{\#\ker(A_b)}{b^{\,n-m}},
\qquad
\mathcal H_A(s):=\sum_{b\ge 1}\frac{h_A(b)}{b^s},
\qquad
\Delta_A:=\prod_{i=1}^{m}d_i.
\]
对每个素数 \(p\)，记
\[
e_i(p):=v_p(d_i),
\qquad
e_p:=\max_i e_i(p),
\qquad
s_p(k):=\sum_{i=1}^{m}\min(k,e_i(p)).
\]
则有精确 Euler 分解
\[
\mathcal H_A(s)
=
\zeta(s)\prod_{p\mid \Delta_A}Q_{A,p}(p^{-s}),
\]
其中每个局部修正项都是正系数多项式
\[
Q_{A,p}(X)
=
1+\sum_{k=1}^{e_p}\bigl(p^{s_p(k)}-p^{s_p(k-1)}\bigr)X^k
\in 1+X\ZZ_{>0}[X].
\]
特别地，\(\mathcal H_A(s)\) 可亚纯延拓到整个复平面，且唯一极点是 \(s=1\) 处的单极点。
\leanverified{paper\_conclusion\_smith\_normalized\_kernel\_positive\_finite\_euler\_correction}
\end{theorem}
\begin{proof}
由推论 \ref{cor:killo-kernel-size-general-modulus}，
\[
h_A(b)=\prod_{i=1}^{m}\gcd(d_i,b),
\]
故 \(h_A\) 为乘法函数，从而
\[
\mathcal H_A(s)=\prod_p L_{A,p}(s),
\qquad
L_{A,p}(s):=\sum_{k\ge 0}\frac{h_A(p^k)}{p^{ks}}.
\]
再由定理 \ref{thm:killo-smith-loss-spectrum}，
\[
h_A(p^k)=p^{s_p(k)}.
\]
若 \(p\nmid \Delta_A\)，则全部 \(e_i(p)=0\)，故 \(s_p(k)=0\) 对一切 \(k\) 成立，从而
\[
L_{A,p}(s)=\sum_{k\ge 0}p^{-ks}=\frac{1}{1-p^{-s}}.
\]

以下设 \(p\mid\Delta_A\)。令 \(a_k:=p^{s_p(k)}\)。当 \(k\ge e_p\) 时，\(\min(k,e_i(p))=e_i(p)\) 对全部 \(i\) 成立，故 \((a_k)\) 在 \(k\ge e_p\) 后稳定。于是
\[
(1-X)\sum_{k\ge 0}a_kX^k
=
a_0+\sum_{k\ge 1}(a_k-a_{k-1})X^k
=
1+\sum_{k=1}^{e_p}\bigl(p^{s_p(k)}-p^{s_p(k-1)}\bigr)X^k,
\]
即
\[
L_{A,p}(s)=\frac{Q_{A,p}(p^{-s})}{1-p^{-s}}.
\]
把全部素数上的局部因子相乘，得到
\[
\mathcal H_A(s)
=
\prod_{p\nmid\Delta_A}\frac{1}{1-p^{-s}}
\prod_{p\mid\Delta_A}\frac{Q_{A,p}(p^{-s})}{1-p^{-s}}
=
\zeta(s)\prod_{p\mid\Delta_A}Q_{A,p}(p^{-s}).
\]
对 \(1\le k\le e_p\)，有
\[
s_p(k)-s_p(k-1)=\#\{i:e_i(p)\ge k\}\ge 1,
\]
故
\[
p^{s_p(k)}-p^{s_p(k-1)}\in\ZZ_{>0}.
\]
因此每个 \(Q_{A,p}\) 都具有正整数系数。有限乘积 \(\prod_{p\mid\Delta_A}Q_{A,p}(p^{-s})\) 是整函数，故 \(\mathcal H_A(s)\) 的唯一极点只能来自 \(\zeta(s)\) 在 \(s=1\) 的单极点。证毕。
\end{proof}

\begin{corollary}[规范化模核计数的平均阶主项与 Smith 扭结检测常数]\label{cor:conclusion-smith-normalized-kernel-average-order}
沿用定理 \ref{thm:conclusion-smith-normalized-kernel-positive-finite-euler-correction} 的记号，记
\[
R_A:=\prod_{p\mid\Delta_A}Q_{A,p}(p^{-1}).
\]
则有：
\begin{enumerate}
  \item \(R_A>0\)，并且
  \[
  R_A=1
  \iff
  \Delta_A=1
  \iff
  d_1=\cdots=d_m=1.
  \]
  \item 规范化模核计数满足显式平均阶
  \[
  \sum_{b\le x}\frac{\#\ker(A_b)}{b^{\,n-m}}
  =
  \sum_{b\le x}h_A(b)
  \sim
  R_A\,x
  \qquad (x\to\infty).
  \]
\end{enumerate}
\leanverified{paper\_conclusion\_smith\_normalized\_kernel\_average\_order}
\end{corollary}
\begin{proof}
由定理 \ref{thm:conclusion-smith-normalized-kernel-positive-finite-euler-correction}，每个 \(Q_{A,p}(X)\) 具有正整数系数，因此
\[
Q_{A,p}(1/p)>0,
\]
从而 \(R_A>0\)。

若 \(\Delta_A=1\)，则空乘积给出 \(R_A=1\)。反之若 \(\Delta_A>1\)，则存在某个 \(p\mid\Delta_A\)。此时 \(e_p\ge 1\)，而
\[
Q_{A,p}(1/p)
=
1+\sum_{k=1}^{e_p}\bigl(p^{s_p(k)}-p^{s_p(k-1)}\bigr)p^{-k}
>1.
\]
其余因子都为正，因此 \(R_A>1\)。故
\[
R_A=1\iff \Delta_A=1.
\]
又 \(\Delta_A=\prod_i d_i\) 且每个 \(d_i\ge 1\)，故 \(\Delta_A=1\) 等价于全部 \(d_i=1\)。

再由定理 \ref{thm:conclusion-smith-normalized-kernel-positive-finite-euler-correction}，
\[
\mathcal H_A(s)=\zeta(s)\prod_{p\mid\Delta_A}Q_{A,p}(p^{-s}),
\]
其中乘积因子在 \(\Re s\ge 1\) 上全纯，故 \(\mathcal H_A(s)\) 在 \(s=1\) 的留数正是 \(R_A\)。由于 \(h_A(b)\ge 0\) 对全部 \(b\) 成立，应用 Wiener--Ikehara 定理即可得到
\[
\sum_{b\le x}h_A(b)\sim R_Ax.
\]
代回 \(h_A(b)=\#\ker(A_b)/b^{\,n-m}\) 即得结论。证毕。
\end{proof}


\begin{theorem}[Smith 亏损轮廓族的完全分类]\label{thm:conclusion-smith-loss-profile-family-classification}
固定整数 \(m\le n\)。对一族函数
\[
f_p:\NN_{\ge 0}\to \NN
\qquad (p\in\mathbb P),
\]
考虑其是否能表示为某个满行秩整数矩阵 \(A\in M_{m\times n}(\ZZ)\) 的 Smith 亏损轮廓
\[
f_{A,p}(k):=\log_p\#\ker(A\bmod p^k)-k(n-m).
\]
则这当且仅当以下四条件同时成立：
\begin{enumerate}
  \item \(f_p(0)=0\)；
  \item
  \[
  \Delta_p(k):=f_p(k)-f_p(k-1)
  \]
  是取值于 \(\{0,1,\dots,m\}\) 的非增序列；
  \item \(\Delta_p(k)=0\) 当 \(k\gg 1\)；
  \item \(f_p\not\equiv 0\) 的素数 \(p\) 只有有限多个。
\end{enumerate}
\end{theorem}
\begin{proof}
必要性来自定理 \ref{thm:conclusion-smith-bounded-rank-loss-profile-characterization} 逐素数的刻画：若
\[
f_{A,p}(k)=\sum_{i=1}^m \min(k,e_i(p)),
\]
则
\[
\Delta_p(k)=\#\{i:\ e_i(p)\ge k\}
\]
自然是取值于 \(\{0,\dots,m\}\) 的非增序列，并对充分大 \(k\) 消失。又只有有限多个素数整除某个 Smith 不变量 \(d_i\)，故非零轮廓只出现在有限多个素数上。

反之，若给定满足四条件的一族 \((f_p)\)，则对每个 \(p\) 可由定理 \ref{thm:conclusion-smith-bounded-rank-loss-profile-characterization} 恢复一组降序估值
\[
e_1^\downarrow(p)\ge \cdots \ge e_m^\downarrow(p)\ge 0
\]
使
\[
f_p(k)=\sum_{i=1}^m \min(k,e_i^\downarrow(p)).
\]
把它们改写为升序列
\[
e_i^\uparrow(p):=e_{m+1-i}^\downarrow(p)
\qquad (1\le i\le m),
\]
并定义
\[
d_i:=\prod_p p^{e_i^\uparrow(p)}.
\]
由于仅有有限多个 \(p\) 使 \(f_p\not\equiv0\)，故上式是有限乘积；又 \(e_i^\uparrow(p)\le e_{i+1}^\uparrow(p)\) 对每个 \(p\) 成立，故
\[
d_1\mid d_2\mid \cdots \mid d_m.
\]
取 Smith 对角矩阵
\[
A=\bigl[\operatorname{diag}(d_1,\dots,d_m)\ \ 0_{m\times(n-m)}\bigr],
\]
则其 \(p\)-主估值正是 \((e_i^\uparrow(p))\)，于是对应亏损轮廓恰为给定的 \(f_p\)。证毕。
\end{proof}

\begin{corollary}[Smith 亏损轮廓的 Euler 原子分解]\label{cor:conclusion-smith-loss-euler-atomic-decomposition}
\leanverified{paper\_conclusion\_smith\_loss\_euler\_atomic\_decomposition}
设 \(A\) 为满行秩整数矩阵。对每个素数 \(p\) 与正整数 \(e\)，定义原子轮廓
\[
\mathfrak a_{p,e}(k):=\min(k,e).
\]
则 \(A\) 的全部 Smith 亏损轮廓唯一分解为
\[
f_{A,p}(k)=\sum_{i=1}^m \mathfrak a_{p,e_i(p)}(k),
\]
并且若 \(A,B\) 为满行秩整数矩阵，则
\[
f_{A\oplus B,p}=f_{A,p}+f_{B,p}.
\]
因此，所有 Smith 亏损轮廓在直和运算下构成由原子 \(\mathfrak a_{p,e}\) 自由生成的交换幺半群。
\end{corollary}
\begin{proof}
由定理 \ref{thm:conclusion-smith-loss-profile-family-classification} 与定理 \ref{thm:conclusion-smith-bounded-rank-loss-profile-characterization}，
\[
f_{A,p}(k)=\sum_{i=1}^m \min(k,e_i(p))
\]
正是原子轮廓的逐项求和。唯一性来自差分剖面
\[
\Delta_p(k)=f_{A,p}(k)-f_{A,p}(k-1)=\#\{i:\ e_i(p)\ge k\},
\]
因为它可唯一恢复估值多重集 \(\{e_i(p)\}\)。对直和 \(A\oplus B\) 而言，Smith 不变量多重集合按并相加，故对应亏损轮廓逐点相加。证毕。
\end{proof}

\begin{proposition}[坏素集合的亏损谱刻画]\label{prop:conclusion-smith-bad-prime-loss-support}
\leanverified{paper\_conclusion\_smith\_bad\_prime\_loss\_support}
对任意满行秩整数矩阵 \(A\in M_{m\times n}(\ZZ)\)，定义其坏素集合
\[
\mathrm{Bad}(A):=\{p\in\mathbb P:\exists\, i,\ v_p(d_i)>0\},
\]
其中 \(d_1\mid\cdots\mid d_m\) 为 \(A\) 的 Smith 不变量。则
\[
p\in \mathrm{Bad}(A)
\iff
f_{A,p}\not\equiv 0
\iff
\exists k\ge 1,\ \#\ker(A\bmod p^k)>p^{k(n-m)}.
\]
因而，矩阵的全部算术坏素恰由局部亏损轮廓的支撑集精确恢复。
\end{proposition}
\begin{proof}
若 \(p\notin \mathrm{Bad}(A)\)，则所有 \(e_i(p)=v_p(d_i)=0\)，从而
\[
f_{A,p}(k)=\sum_{i=1}^m \min(k,0)=0
\qquad(\forall k).
\]
反之，若某个 \(e_i(p)>0\)，则
\[
f_{A,p}(1)\ge 1,
\]
故 \(f_{A,p}\not\equiv0\)。这证明了前两个条件的等价。

另一方面，由定理 \ref{thm:conclusion-smith-kernel-dirichlet-euler-rational-local-factors} 的局部核公式，
\[
\#\ker(A\bmod p^k)=p^{k(n-m)+f_{A,p}(k)}.
\]
因此
\[
f_{A,p}(k)>0
\iff
\#\ker(A\bmod p^k)>p^{k(n-m)}
\]
对某个 \(k\) 成立。于是三者等价。证毕。
\end{proof}

\begin{theorem}[Smith \texorpdfstring{$p$}{p}-主谱的 Ramanujan 素数幂影公式]\label{thm:conclusion-smith-pprimary-ramanujan-shadow-formula}
设 \(A\in M_{m\times n}(\ZZ)\) 满足 \(\operatorname{rank}_{\QQ}(A)=m\)，其 Smith 不变量记为
\[
d_1\mid d_2\mid\cdots\mid d_m.
\]
对素数 \(p\)，记
\[
e_i(p):=v_p(d_i),
\qquad
\Delta_p(k):=\#\{\,i:e_i(p)\ge k\,\}
\quad (k\ge 1),
\qquad
\Delta_p(0):=m,
\]
并定义 prime-power Ramanujan 影
\[
\mathcal R_{A,p}(k):=\sum_{i=1}^{m}c_{p^k}(d_i)
\qquad (k\ge 1).
\]
则对任意 \(k\ge 1\)，成立精确恒等式
\[
\mathcal R_{A,p}(k)=p^k\Delta_p(k)-p^{k-1}\Delta_p(k-1).
\]
因此，单个素数塔上的 Ramanujan 影精确记录了 \(p\)-主 Smith 图样的列高差分。
\leanverified{shadow\_formula\_algebraic}
\leanverified{shadow\_formula\_telescoping\_K3}
\leanverified{paper\_conclusion\_smith\_ramanujan\_shadow\_formula}
\end{theorem}
\begin{proof}
对素数幂导子，Ramanujan 和满足
\[
c_{p^k}(n)=
\begin{cases}
0, & v_p(n)\le k-2,\\[2pt]
-p^{k-1}, & v_p(n)=k-1,\\[2pt]
p^{k-1}(p-1), & v_p(n)\ge k.
\end{cases}
\]
因此
\[
\mathcal R_{A,p}(k)
=
p^{k-1}(p-1)\Delta_p(k)-p^{k-1}\bigl(\Delta_p(k-1)-\Delta_p(k)\bigr),
\]
整理即得
\[
\mathcal R_{A,p}(k)=p^k\Delta_p(k)-p^{k-1}\Delta_p(k-1).
\]
证毕。
\end{proof}

\begin{theorem}[prime-power Ramanujan 影的完全反演与实现判别]\label{thm:conclusion-smith-pprimary-ramanujan-shadow-inversion}
固定素数 \(p\) 与整数 \(m\ge 1\)。对任意数列 \((R_k)_{k\ge 1}\)，定义
\[
\Delta_k:=p^{-k}\left(m+\sum_{j=1}^{k}R_j\right)
\qquad (k\ge 1).
\]
则 \((R_k)_{k\ge 1}\) 可表示成某个满行秩整数矩阵的 \(p\)-主 Ramanujan 影
\[
R_k=\mathcal R_{A,p}(k)
\qquad (k\ge 1)
\]
当且仅当 \((\Delta_k)_{k\ge 1}\) 是一个非负整数、单调不增并最终为零的序列。并且在该情形下，
\[
\#\{\,i:e_i(p)=k\,\}=\Delta_k-\Delta_{k+1}.
\]
\leanverified{paper_conclusion_smith_ramanujan_shadow_inversion_seeds}
\leanverified{paper_conclusion_smith_ramanujan_shadow_inversion_package}
\leanverified{paper\_conclusion\_smith\_pprimary\_ramanujan\_shadow\_inversion}
\end{theorem}
\begin{proof}
若 \(R_k=\mathcal R_{A,p}(k)\)，则由定理 \ref{thm:conclusion-smith-pprimary-ramanujan-shadow-formula}，
\[
\mathcal R_{A,p}(k)=p^k\Delta_p(k)-p^{k-1}\Delta_p(k-1).
\]
令 \(E_k:=p^k\Delta_p(k)\)，则
\[
\mathcal R_{A,p}(k)=E_k-E_{k-1},
\qquad
E_0=\Delta_p(0)=m.
\]
望远镜求和即得
\[
E_k=m+\sum_{j=1}^{k}\mathcal R_{A,p}(j),
\]
亦即
\[
\Delta_p(k)=p^{-k}\left(m+\sum_{j=1}^{k}\mathcal R_{A,p}(j)\right).
\]
因此 \((\Delta_k)\) 必为非负整数、单调不增且最终为零。再由列高差分公式，
\[
\#\{\,i:e_i(p)=k\,\}=\Delta_p(k)-\Delta_p(k+1).
\]

反之，若 \((\Delta_k)\) 满足上述三个条件，则
\[
m_k:=\Delta_k-\Delta_{k+1}\in \ZZ_{\ge 0},
\qquad
\sum_{k\ge 0}m_k=m.
\]
取恰有 \(m_k\) 个 \(p\)-主赋值等于 \(k\) 的多重集 \(\{e_i(p)\}_{i=1}^{m}\)，再取对应 Smith 对角矩阵即可实现该多重集。按定理 \ref{thm:conclusion-smith-pprimary-ramanujan-shadow-formula} 反算，即得 \(R_k=\mathcal R_{A,p}(k)\)。证毕。
\end{proof}

\begin{theorem}[局部 zeta 因子与 Ramanujan 影的双向有限差分同构]\label{thm:conclusion-smith-local-zeta-ramanujan-shadow-difference-equivalence}
在上述记号下，定义
\[
f_p(-1):=-m,
\qquad
f_p(0):=0,
\qquad
f_p(k):=\sum_{i=1}^{m}\min(k,e_i(p))
\quad (k\ge 1).
\]
则对一切 \(k\ge 1\)，成立二阶有限差分恒等式
\[
\mathcal R_{A,p}(k)
=
p^{k-1}\Bigl(
p\,f_p(k)-(p+1)f_p(k-1)+f_p(k-2)
\Bigr).
\]
反过来，
\[
f_p(k)
=
\sum_{j=1}^{k}\Delta_p(j)
=
\sum_{j=1}^{k} p^{-j}\left(m+\sum_{\ell=1}^{j}\mathcal R_{A,p}(\ell)\right).
\]
因而，局部因子
\[
Z_{A,p}(s)=\sum_{k\ge 0}p^{-k(s-n+m)+f_p(k)}
\]
与 Ramanujan 影序列 \((\mathcal R_{A,p}(k))_{k\ge 1}\) 彼此唯一决定。
\leanverified{paper\_conclusion\_smith\_local\_zeta\_ramanujan\_shadow\_difference\_equivalence}
\end{theorem}
\begin{proof}
由定义，
\[
f_p(k)-f_p(k-1)=\Delta_p(k)
\qquad (k\ge 1).
\]
把它代入定理 \ref{thm:conclusion-smith-pprimary-ramanujan-shadow-formula} 的
\[
\mathcal R_{A,p}(k)=p^k\Delta_p(k)-p^{k-1}\Delta_p(k-1),
\]
便得
\[
\mathcal R_{A,p}(k)
=
p^{k-1}\Bigl(
p(f_p(k)-f_p(k-1))-(f_p(k-1)-f_p(k-2))
\Bigr),
\]
即所述二阶差分公式。

另一方面，由上一条定理，
\[
\Delta_p(j)=p^{-j}\left(m+\sum_{\ell=1}^{j}\mathcal R_{A,p}(\ell)\right),
\]
对 \(j=1,\dots,k\) 求和即得 \(f_p(k)\) 的反演公式。最后把 \((f_p(k))\) 代入局部 zeta 因子闭式，便知 \((\mathcal R_{A,p}(k))\) 与 \(Z_{A,p}\) 彼此唯一决定。证毕。
\end{proof}

\begin{proposition}[Smith 光谱的 Ramanujan--Dirichlet 阴影]\label{prop:conclusion-smith-global-ramanujan-dirichlet-shadow}
\leanverified{paper\_conclusion\_smith\_global\_ramanujan\_dirichlet\_shadow}
定义全局 Ramanujan 影
\[
\mathcal R_A(q):=\sum_{i=1}^{m}c_q(d_i),
\]
以及其 Dirichlet 级数
\[
\mathfrak R_A(s):=\sum_{q\ge 1}\frac{\mathcal R_A(q)}{q^s}.
\]
则当 \(\Re s>1\) 时，
\[
\mathfrak R_A(s)=\frac{1}{\zeta(s)}\sum_{i=1}^{m}\sigma_{1-s}(d_i).
\]
\leanverified{paper\_conclusion\_smith\_global\_ramanujan\_dirichlet\_shadow}
\end{proposition}
\begin{proof}
经典恒等式给出
\[
\sum_{q\ge 1}\frac{c_q(n)}{q^s}=\frac{\sigma_{1-s}(n)}{\zeta(s)}
\qquad (\Re s>1).
\]
对 \(n=d_i\) 求和并交换有限求和次序，即得
\[
\mathfrak R_A(s)
=
\sum_{i=1}^{m}\sum_{q\ge 1}\frac{c_q(d_i)}{q^s}
=
\frac{1}{\zeta(s)}\sum_{i=1}^{m}\sigma_{1-s}(d_i).
\]
证毕。
\end{proof}

\begin{theorem}[有限 Ramanujan 审计的不完备性]\label{thm:conclusion-smith-truncated-ramanujan-audit-incompleteness}
固定有限素数集 \(S\subset \PP\) 与截断高度 \(K\ge 1\)。定义截断 Ramanujan 审计数据
\[
\mathfrak R_{S,K}(A):=\{\mathcal R_{A,p}(k):p\in S,\ 1\le k\le K\}.
\]
则 \(\mathfrak R_{S,K}\) 不是完整 Smith 光谱的完备不变量；更强地说，每一个固定的 \(\mathfrak R_{S,K}\)-等价类都包含无穷多个两两不同的 Smith 光谱。
\leanverified{paper\_conclusion\_smith\_truncated\_ramanujan\_audit\_incompleteness}
\end{theorem}
\begin{proof}
由定理 \ref{thm:conclusion-truncated-smith-ledger-infinite-indistinguishability}，对任意固定的 \(S\) 与 \(K\)，都存在无穷多个满行秩整数矩阵 \(A^{(r)}\) 使得
\[
K_p(k;A^{(r)})=K_p(k;A)
\qquad
(p\in S,\ 1\le k\le K),
\]
而它们的完整 Smith 不变量系两两不同。
由局部核公式，
\[
K_p(k;A)=p^{k(n-m)+f_{A,p}(k)},
\]
故这些矩阵在同一有限窗口内具有完全相同的 \(f_{A,p}(k)\)。再由定理 \ref{thm:conclusion-smith-local-zeta-ramanujan-shadow-difference-equivalence}，它们在该窗口内的 \(\mathcal R_{A,p}(k)\) 亦完全相同。于是 \(\mathfrak R_{S,K}\) 不能唯一决定完整 Smith 光谱。证毕。
\end{proof}

\begin{theorem}[prime-power Ramanujan 审计的精确 cutoff 完备性]\label{thm:conclusion-smith-pprimary-ramanujan-cutoff-completeness}
\leanverified{paper\_conclusion\_smith\_pprimary\_ramanujan\_cutoff\_completeness}
对固定素数 \(p\)，记
\[
E_p^{\max}:=\max_{1\le i\le m}e_i(p).
\]
则截断序列
\[
\bigl(\mathcal R_{A,p}(1),\dots,\mathcal R_{A,p}(K)\bigr)
\]
唯一决定 \(A\) 的 \(p\)-主 Smith 光谱，当且仅当
\[
K\ge E_p^{\max}.
\]
\end{theorem}
\begin{proof}
若 \(K\ge E_p^{\max}\)，则对所有 \(k>K\) 有 \(\Delta_p(k)=0\)。由定理 \ref{thm:conclusion-smith-pprimary-ramanujan-shadow-inversion}，
\[
\Delta_p(k)=p^{-k}\left(m+\sum_{j=1}^{k}\mathcal R_{A,p}(j)\right)
\]
已由截断数据完全确定；而当 \(k>K\) 时自动为零。因此全部 \(\Delta_p(k)\) 都可恢复，进而
\[
\#\{\,i:e_i(p)=k\,\}=\Delta_p(k)-\Delta_p(k+1)
\]
给出完整 \(p\)-主光谱。

反之，若 \(K<E_p^{\max}\)，则存在至少一个 \(i\) 使 \(e_i(p)>K\)。把该指数从 \(E\) 改成任意更大的 \(E'>E\)，则对所有 \(1\le k\le K\)，\(\Delta_p(k)\) 保持不变，从而按定理 \ref{thm:conclusion-smith-pprimary-ramanujan-shadow-formula}，
\[
\mathcal R_{A,p}(k)
\qquad
(1\le k\le K)
\]
亦保持不变；然而 \(p\)-主 Smith 光谱已经改变。故当 \(K<E_p^{\max}\) 时，截断序列不可能完备。证毕。
\end{proof}

\begin{conjecture}[Ramanujan prime-power 影的初对象性]\label{conj:conclusion-smith-ramanujan-shadow-initiality}
在一切满足以下三条的 \(p\)-局部不变量族中：
\begin{enumerate}
\normalfont
  \item 仅依赖于 \(\{e_i(p)\}_{i=1}^{m}\)；
  \item 对 block sum 可加，或等价地可经有限差分递推恢复；
  \item 对有限 truncation 具有精确 cutoff 完备门槛，
\end{enumerate}
prime-power Ramanujan 影族
\[
\{\mathcal R_{A,p}(k)\}_{k\ge 1}
\]
应当是初对象：任何其他此类完备族，都应经由一个上三角、有限差分型的有理变换从 \(\{\mathcal R_{A,p}(k)\}\) 导出。

其根据是刚性的：定理 \ref{thm:conclusion-smith-pprimary-ramanujan-shadow-formula} 表明它与列高函数 \(\Delta_p(k)\) 一阶望远镜等价，定理 \ref{thm:conclusion-smith-local-zeta-ramanujan-shadow-difference-equivalence} 表明它与局部 zeta 因子二阶有限差分等价，而定理 \ref{thm:conclusion-smith-pprimary-ramanujan-cutoff-completeness} 又给出恰好由 \(E_p^{\max}\) 控制的精确 cutoff 门槛。因此，若还存在另一族同等完备的 \(p\)-局部坐标，它应当只是 Ramanujan 影的结构性重写，而不是新的原子层。
\end{conjecture}


\noindent
由此，Smith 亏损谱在结论层完成了从“函数轮廓”到“有限证书”的严格压缩：固定行秩以后，亏损轮廓与长度恰为 \(m\) 的估值分拆、其 Young 对偶列高、以及坏素数上的局部形式幂级数彼此完全等价。换言之，局部亏损函数、局部 Euler 因子与全局 Smith 不变量之间不再存在任何信息损失；真正需要跟踪的无限对象，已经塌缩为有限个坏素数上的有限深度赋值数据。

\endinput
